\chapter{Практика}
\label{ch:chap2}


\section{Описание знакового кодера}

Задача знакового кодера – преобразовать входное символьное сообщение в набор двоичных символов (бит). 
На вход функции поступает текстовое сообщение длительностью от 30 до 100 символов. 
Длина сообщения задается вариантом. Алфавит сообщения состоит из 64 символов: латинский (маленькие и заглавные буквы, 52 символа), 
цифры (0…9), пробел и точка. Требуется преобразовать текстовое сообщение в битовое. 
Для этого нужно выбрать разрядность кодового слова (количество бит на 1 символ сообщения), а каждый символ должен быть закодирован индивидуальным кодом. 
На выходе функции должно быть битовое сообщение.
Запрещено использовать встроенную в функцию преобразования
формата данных.

\begin{figure}[H]
    \centering
    \includegraphics{image1.png}
    \caption{Блок схема знакового кодера/декодера}
\end{figure}

\section{Описание ASCII таблицы}

На рисунке ниже показан пример символьного кодирования слова TUSUR
с помощью алфавита ASCII, согласно которому каждый символ
кодируется семизначным кодовым словом.

\noindent
\begin{minipage}{0.5\textwidth}
\centering
\includegraphics[width=\linewidth]{image3.png}
\end{minipage}
\hfill
\begin{minipage}{0.5\textwidth}
\centering
\includegraphics[width=\linewidth]{image2.png}
\end{minipage}


\newpage

\noindent
[\dotfill]

\section{Описание входных выходных переменных кодера}

\textbf{Входные переменные:}  вектор с текстовым сообщением (битами)

\textbf{Выходные переменные:} вектор с битами

\textbf{Параметры:} нет


\section{Описание знакового декодера}
Знаковый декодер — это обратная функция знакового кодера, выполняющая инверсную операцию. 
На его вход поступает битовое сообщение произвольной длины. Требуется преобразовать это битовое сообщение в текстовое, 
используя такую же знаковую кодировку, как в функции
знакового кодирования. На выходе функции должно быть такое же текстовое сообщение, как и на входе знакового кодера.

\noindent
[\dotfill]

\section{Описание входных выходных переменных декодера}

\textbf{Входные переменные:}  вектор с битами

\textbf{Выходные переменные:} вектор с текстовым сообщением

\textbf{Параметры:} нет



% \section{Задания}

% Сгенерируем СП с нормальным и равномерным распределением. Также сгенерируем сигнал в виде
% синусоиды и наложим на сигнал шум:

% \begin{figure}[H]
%     \centering
%     \includegraphics{10.png}
%     \caption{Сигнал и разные шумы}
% \end{figure}

% \begin{figure}[H]
%     \centering
%     \includegraphics{20.png}
%     \caption{Зашумленный сигнал}
% \end{figure}

% Теперь вычислим АКФ:

% \begin{figure}[H]
%     \centering
%     \includegraphics{30.png}
%     \caption{АКФ сигнала с шумом}
% \end{figure}

% Можно увидеть, что шум исчез, а благодаря периодичности сигнала
% АКФ похожа на исходный сигнал.\\

% Повысим дисперсию нормального шума. Теперь шум по амплитуде больше сигнала.
% В таком случае АКФ получится не такой гладкой, но она всё равно напоминает исходный сигнал


% \begin{figure}[H]
%     \centering
%     \includegraphics{40.png}
%     \caption{АКФ сигнала с сильным шумом}
% \end{figure}

% Случайный процесс состоит из суммы быстрых и медлленных изменений. Эти компоненты можно выделить с помощью ФНЧ 
% или ФВЧ

% \begin{figure}[H]
%     \centering
%     \includegraphics{50.png}
%     \caption{ИХ фильтра}
% \end{figure}

% \begin{figure}[H]
%     \centering
%     \includegraphics{60.png}
%     \caption{Быстрая и медленная компонента}
% \end{figure}

% Теперь посчитаем АКФ для этих процессов

% \begin{figure}[H]
%     \centering
%     \includegraphics{70.png}
%     \caption{АКФ для медленной и быстрой компоненты}
% \end{figure}

% \begin{figure}[H]
%     \centering
%     \includegraphics{80.png}
%     \caption{АКФ для медленной и быстрой компоненты}
% \end{figure}

% Если приблизить АКФ, то можно заметить, что интервал корреляции
% быстрой компоненты меньше, чем интервал корреляции медленной компоненты.
% Это говорит о том, что по интервалу корреляции мы можем судить о поведении
% процесса, а это очень важно, потому что по числовым характеристикам
% процесса мы не можем судить о его поведении, а АКФ как раз это может показать.







\endinput 